\chapter{A worked example: from a raw SAM to results}\label{chap:example}

This chapter takes an unbalanced SAM through every step of the workflow: a
template of roles, aggregation and balancing, mapping, elasticities from the
library, calibration and checks, three policy simulations, sensitivity
analysis, a dynamic baseline and scenario, the link to a household survey, and
saving the model. The SAM and the household survey are SYNTHETIC: they are
constructed for testing and describe no country. The elasticities are
published values for Togo, applied to the synthetic SAM only to show the
commands, so the results illustrate the workflow and are not findings about
any country. The commands are those of the do-file
\cmd{raw_sam_to_results.do} in the \cmd{examples} folder of the package's
repository.

\section{The raw SAM}

\stlog{we_raw}
The SAM has 17 accounts; its largest difference between an account's row and
column totals is 12.1. It has two household accounts, households (\cmd{hhd})
and non-profit institutions (\cmd{npi}), and an enterprise (\cmd{ent}).

\section{Roles, aggregation and balancing}

\manref{cge_template} writes a template with one row per account, to be filled
in, with a real SAM usually in Excel:
\stlog{we_template}
The do-file fills in the roles and puts \cmd{hhd} and \cmd{npi} in one
household group, \cmd{hh}:
\stlog{we_roles}
The last three columns show where the SAM does not balance. \manref{cge_aggregate}
merges the two household accounts, and \manref{cge_balance} balances the SAM
by cross entropy, with the income of the household group taken as known
(\secref{sec:ce}):
\stlog{we_balance}
The method converges in four iterations. Every account now balances, 61 of
the 63 non-zero cells have changed, and the SAM total has risen from 1,860.5 to
1,943.8. The known household income, 290, is above the row total of the merged
household accounts in the raw SAM, about 257, and the largest changes are in
the flows that pay and spend it: the labour income of households, for
example, rises by 10.5 per cent. These are the changes to report with any result. The
same template then maps the balanced SAM; the options it implies are
displayed, so that the mapping can also be written as a command:
\stlog{we_map}

\section{Elasticities and calibration}

\manref{cge_elasticities} returns values from the library as options for
\manref{cge_calibrate}: the aggregate Armington and CET elasticities estimated
for Togo, the capital--labour elasticity estimated for Africa, income
elasticities for food (mapped to agriculture) and for other goods and
services, and the Frisch parameter (chapter~\ref{chap:elasticities}):
\stlog{we_elas}
With income elasticities and a Frisch parameter, household demand is a linear
expenditure system:
\stlog{we_calibrate}
\manref{cge_parameters} shows the implied rates: import tariffs of 10.6 per
cent on agriculture and 9.8 per cent on manufactures, sales taxes of 5.7 to 7.7
per cent, and direct tax rates of 6.6 per cent for the household group and 9.9
per cent for the enterprise:
\stlog{we_parameters}

\section{Three policy simulations}

The do-file removes every import tariff under the default closure and with
direct tax rates replacing the lost revenue, and then raises government
purchases of services by 20 per cent under the default closure:
\stlog{we_sims}
\stlog{we_report}
The closure matters as much as the shock. With government savings flexible,
the lost tariff revenue (8.2, or 11 per cent of revenue) widens the deficit;
with savings-driven investment, investment falls by 11.9 per cent and real
consumption rises by 2.5 per cent. With direct taxes replacing the revenue,
the direct tax take rises by 29.8 per cent, real consumption falls by 0.35 per
cent and investment rises by 4.0 per cent. The exchange rate depreciates by 4.6
to 4.7 per cent in both cases, and exports rise. More government purchases of
services lower investment by 16.6 per cent under the same closure.
\manref{cge_report} with \cmd{saving()} writes every table to Excel.

\section{Sensitivity}

\manref{cge_sensitivity} recalibrates the model for 50 draws of the
elasticities, each within plus or minus 50 per cent of the central values,
sector by sector, and runs the tariff removal:
\stlog{we_sens}
All 50 draws converge. The welfare result barely moves (equivalent variation
2.48 per cent; 5th to 95th percentile 2.48 to 2.49), nor does investment
($-12.2$ to $-11.6$ per cent), because both are set mainly by the closure.
Trade responds more: exports rise by 0.9 to 1.8 per cent, and the exchange rate
depreciates by 4.0 to 5.2 per cent.

\section{Dynamics}

A baseline grows real value added by 4 per cent a period, with capital that
accumulates from investment and is specific to each activity; the efficiency
factors that achieve it fall from 2.4 to 2.0 per cent a period. The scenario
removes tariffs from period 1 with the baseline's efficiency path:
\stlog{we_dynamic}
\manref{cge_dynresults} compares the scenario with the baseline:
\stlog{we_dynresults}
Because the static closure lowers investment, capital grows more slowly: by
period 5, output is 2.2 per cent below the baseline in manufacturing, 1.4 per
cent in agriculture and 0.8 per cent in services. The static gain in
consumption has a dynamic cost, which a static run does not show.

\section{Poverty and inequality}

The SYNTHETIC survey has 1,000 households with wages, profits and government
transfers, and budget shares for food, goods and services. \manref{cge_microsim}
passes the simulation with direct taxes replacing the tariff revenue to it:
\stlog[\scriptsize]{we_microsim}
Real consumption per person rises by 0.64 per cent; the headcount ratio falls
from 23.09 to 22.97 per cent, and the Gini coefficient barely changes. The
households' gain differs from the group's equivalent variation ($-0.36$ per
cent) because they keep their own income composition (here mostly wages, with
no transfers from enterprises) and budget shares, and because real
consumption is not a welfare measure. A real survey would map every income
source of the households.

\section{Saving the model}

\stlog{we_save}
\manref{cge_use} reloads the model with its 15 simulations, and
\manref{cge_report} reproduces the results without solving again.
